\section{Open problem}
\label{sec:open}

Corollary~\ref{cor:impossible} shows that no reset staying inside the small-error
$\varepsilon$-linear class can improve the worst-case cross-talk scale beyond $d^{-1/2}$ when
$F\gg d$, while Theorem~\ref{thm:randomsupport} shows thresholded recovery survives to
$s=O(d/\log d)$ at quadratic load. Any interpolation beyond the H\"{a}nni template must
therefore leave the small-error linear class. This makes one question concrete.

\begin{conjecture}[Generic robust threshold reset]
\label{conj:reset}
Fix $s=O(1)$. There are constants $c,C,K>0$ such that for every $d$ and every
$F\le c\,d^2/\log^C d$ there exist a code $\Phi\in\R^{d\times F}$, an input readout
$R_{\mathrm{in}}$, a width-$\Otilde(d)$ reset module $\mathcal R_\Phi:\R^d\to\R^d$ and a
scoring map $S$ such that, for every $s$-sparse $b$ and every $x$ with
$\norm{R_{\mathrm{in}}x-b}_\infty\le K\,d^{1/2}/(F^{1/2}s^{1/4})$, the output
$z=\mathcal R_\Phi(x)$ is composable with the next stage and satisfies
$b_i=1\Rightarrow(Sz)_i\ge3/4$ and $b_i=0\Rightarrow(Sz)_i\le1/4$.
\end{conjecture}

The incoming-error tolerance is the essential clause: a reset theorem for clean inputs
$x=\Phi b$ alone would not support recursive composition. By
Corollary~\ref{cor:impossible} any proof must use an explicitly nonlinear or thresholded
decoding mechanism, as Adler--Shavit do.

% ============================================================================
\section{Conclusion}
\label{sec:conclusion}
% ============================================================================

We proposed the Interface Diagnostic, a computable procedure that decides which decoding
interface an overcomplete code or a trained network maintains, and how far it is from the best
possible one. It rests on a rank--trace floor for unit-diagonal linear readouts, in a
gain-normalised form that bounds the interference-to-gain ratio, and on a separation between the
linear and threshold criteria proved under one and the same sparse-state model. Its mean-square
statistics run in $O(Fd^2)$ time and $O(d^2)$ memory and reach $F\approx10^6$ features on a
laptop CPU.

Two additions make the diagnosis about a code rather than about an ensemble. Computing a code's
own floor in closed form splits any attainment ratio into a geometry term and a decoder term, and
the split is deflationary: under the calibrated pseudoinverse the decoder term is one by
construction, so the ratio a reader is invited to read as a decoder result is nothing of the
kind. The trained $L^4$ network's \EfiveLfourRatioPinvMean{} and the $L^2$ baseline's factor
\EfiveLtwoRatioLsMean{} are statements about leverage profiles. And for support recovery, one
linear programme per feature returns the largest sparsity at which any affine probe can detect
that feature, with a certified margin when it succeeds and a checkable certificate when it fails.
Leverage links the two: low leverage forces the affine frontier down, with no decoder in the
argument, and an explicit code refutes the converse, so the hierarchy holds in one direction
only.

Applied to 180 models, the instrument returns one result that resists summary and two that do not.
The first: given the same representation and the same threshold treatment, the comparison between the
network and a fitted affine probe has no single sign. It is positive under $L^4$ and grows with width
to \ScPrimDiff{} levels in \ScPrimNetWins{} of \ScSeeds{} models at $d=400$; it is exactly zero under
$L^2$, where both decoders fail together; and it is negative by up to
\PrimRandFourHundredDiff{} levels on a code that was never trained. The \PrimTies{} ties out of
\PrimModels{} are therefore not evidence of equivalence --- \PrimLtwoJointFailures{} of them are
joint failures --- and the $L^4$ lead is not evidence of nonlinear expressive power either, since the
probe class contains the network's own decoder (Section~\ref{sec:probes}). What the comparison
identifies is the code, not the decoder. The first positive: what separates
models is the loss, not the architecture --- $L^2$ pushes $R_{\mathrm{geom}}$ to
\SALtwoTwoHundredRgeom{} and collapses its frontier to $1$, while $L^4$ holds
\SALfourTwoHundredRgeom{} and a frontier of \FullLfourTwoHundredKappaMin{}. The second: that
separation is bought by training the code and not the readout, since a frozen random code with a
fully trained readout keeps an untrained code's frontier.

Two things we decline to claim. There is no scaling law here: the threshold-recovery campaign reaches
$d=\EsixDmax$ at about a million features, which establishes that the measurement scales, but four
widths do not distinguish $c\,d/\ln d$ from the alternatives. And proximity to the rank--trace floor
is not evidence of learned structure, because an untrained code already achieves it.

The broader lesson is that capacity questions for overcomplete representations are not settled by
counting features, nor by measuring distance to a bound that random codes already meet. They depend
on which decoding interface a particular code maintains, and on which representation is being read ---
both of which are now measurable, per feature and exactly, for the code in hand.
